\section{The actual normal-image limit, its graph, and its arithmetic observation}
\label{sec:nested-normal-observation}

Keep the finite dagger-stable packet, its entire quotient
\(v_h=g/h\), \(g=2\xi\), and an integer \(k\ge2\). The empty
packet \(h=1\) is included. We retain the original amplitude and coordinates
\[
 a(t)=\frac{v_h(1/2+it)}{\sqrt{2\pi}},\quad
 t_i=u/k+y_i,\quad y_k=-\sum_{i<k}y_i,\quad
 \Psi(u,\mathbf y)=\prod_{i=1}^k a(t_i),\quad S=k/2+iu.
 \tag{NN.1}
\]
The original measures are \(du\,d\mathbf y\), with the unit Jacobian
proved in (FC.2). No mass, phase, residue unit or source coefficient
is replaced. Write \(\mathscr H=L^2(d\mathbf y)\) and
\(\mathscr K=L^2(du;\mathscr H)\).

For each integer \(m\ge0\), let \(\mathcal P_m\) be the complete
original polynomial source of total degree at most \(m\). Its relative
frame consists of every monomial \(\theta_\alpha(\mathbf z)\) of degree
\(\ell_\alpha\le m\) in the \(k-1\) independent variables
\(z_i=s_i-S/k\), \(i<k\), with
\(z_k=-\sum_{i<k}z_i\) retained as the dependent coordinate. Its literal source
basis consists of \(\theta_\alpha S^n\),
\(0\le n\le m-\ell_\alpha\), as in (FC.28). Keep its coefficient
map \(C_m(u)\), frame \(j_m(u)\), orthogonal projection \(\Pi_m(u)\),
connection \(\Gamma_m(u)\), and normal map \(N_m(u)\) of (FC.6).
Thus the actual amplitude, normal and tangential source maps are
\[
 \begin{aligned}
 A_m p&=j_m C_m p=\Psi P_p,\\
 \boldsymbol N_m p&=N_m C_m p=(I-\Pi_m)\partial_u A_m p,\\
 \boldsymbol T_m p&=j_m(C_m'+\Gamma_m C_m)p
                         =\Pi_m\partial_u A_m p.
 \end{aligned}
 \tag{NN.2}
\]
All three maps have values in the same original space \(\mathscr K\).
The last two equalities include the derivative of every original power
of \(S\); its factor is \(inS^{n-1}\), not \(nS^{n-1}\).

Let \(Q_{\ell m}\) be the literal coefficient inclusion for
\(\mathcal P_m\subseteq\mathcal P_\ell\), \(\ell\ge m\), and
\(E_{\ell m}\) its relative-frame inclusion. The coefficient identities
and amplitude identity are
\[
 C_\ell Q_{\ell m}=E_{\ell m}C_m,\qquad
 j_\ell E_{\ell m}=j_m,\qquad A_\ell Q_{\ell m}=A_m.
 \tag{NN.3}
\]
They follow by comparing each literal column \(\theta_\alpha S^n\).
In particular these are inclusions of the original sources, and not
an identification of coefficient spaces of different dimensions.

The weighted-density theorem for the amplitude (NN.1) proves
\(\Pi_m(u)\to I\) strongly for every \(u\), and proves for every
fixed original polynomial \(P\) that
\[
 \begin{aligned}
 r_m(P)&:=(I-\Pi_m)\partial_u(\Psi P)\longrightarrow0
                                      &&\text{in }\mathscr K,\\
 t_m(P)&:=\Pi_m\partial_u(\Psi P)\longrightarrow
                      \partial_u(\Psi P)&&\text{in }\mathscr K.
 \end{aligned}
 \tag{NN.4}
\]
Here and below \(m\ge\deg P\). The full density proof, including its
original exponential estimate, the division across the selected full
zero orders, all fibre zero sets and the invariant-source variant,
is given in the preceding relative-density section. Equivalently,
the first statement of (NN.4) is obtained by applying those actual
fibre projections to the one fixed function
\(\partial_u(\Psi P)\in\mathscr K\), then applying dominated
convergence to the bound
\(\|(I-\Pi_m)\partial_u(\Psi P)\|\le
\|\partial_u(\Psi P)\|\). Thus no unproved estimate about a sequence
of changing source polynomials is used here.

\subsection{Exact contractions of the finite normal images}

Put \(\mathscr R_m^N=\operatorname{ran}\boldsymbol N_m\), and retain
\(H_m^N=\boldsymbol N_m^*\boldsymbol N_m>0\) and
\(L_m^N=(H_m^N)^{-1}\boldsymbol N_m^*\) of (FC.29)--(FC.30).
Their positivity and the exact source inverse hold for the actual
arithmetic amplitude, including \(h=1\). All these source spaces
contain the constant monomial and are nonzero. Define the bounded
orthogonal projection on \(\mathscr K\)
\[
 (\mathbb P_m f)(u)=\Pi_m(u)f(u).
 \tag{NN.5}
\]
The finite-frame entries are smooth, so this is a measurable contraction;
self-adjointness and idempotence hold after integrating their pointwise
identities. The relative inclusions give
\(\mathbb P_\ell\mathbb P_m=\mathbb P_m
=\mathbb P_m\mathbb P_\ell\) for \(\ell\ge m\).

The exact transition from one normal image to the next is
\[
 \begin{aligned}
 V_{\ell m}:\mathscr R_m^N&\longrightarrow\mathscr R_\ell^N,\\
 V_{\ell m}r&=(I-\mathbb P_\ell)r
       =\boldsymbol N_\ell Q_{\ell m}L_m^N r.
 \end{aligned}
 \tag{NN.6}
\]
To prove both its codomain and its second expression, write uniquely
\(r=\boldsymbol N_m p=(I-\mathbb P_m)\partial_u A_m p\).
The nested projection identity and (NN.3) give
\[
 (I-\mathbb P_\ell)r
   =(I-\mathbb P_\ell)\partial_u A_m p
   =\boldsymbol N_\ell Q_{\ell m}p.
\]
Injectivity follows from injectivity of \(\boldsymbol N_\ell\) and
of the literal inclusion. Moreover
\[
 \|V_{\ell m}\|\le1,\quad V_{mm}=I,\quad
 V_{nm}=V_{n\ell}V_{\ell m}\quad(n\ge\ell\ge m).
 \tag{NN.7}
\]
Projection contraction proves the first assertion; the source formula
and \(L_\ell^N\boldsymbol N_\ell=I\) prove the others.
The source and transition square commutes exactly:
\[
 L_\ell^NV_{\ell m}=Q_{\ell m}L_m^N.
 \tag{NN.8}
\]
In particular these changing normal images have a specified algebraic
relation even though their Hilbert subspaces are not asserted to be
nested by the ambient identity map.

Their complete energy difference on the fixed source is
\[
 \begin{aligned}
 r_m(P)&=(\mathbb P_\ell-\mathbb P_m)\partial_u(\Psi P)+r_\ell(P),\\
 \|r_m(P)\|^2-\|r_\ell(P)\|^2
 &=\|(\mathbb P_\ell-\mathbb P_m)\partial_u(\Psi P)\|^2,\\
 \|t_\ell(P)\|^2-\|t_m(P)\|^2
 &=\|(\mathbb P_\ell-\mathbb P_m)\partial_u(\Psi P)\|^2.
 \end{aligned}
 \tag{NN.9}
\]
Indeed \(\mathbb P_\ell-\mathbb P_m\), \(\mathbb P_m\), and
\(I-\mathbb P_\ell\) are mutually orthogonal projections. Apply
their orthogonal decomposition to the unchanged derivative amplitude.
The full nested-frame connection identities in the preceding section
identify this middle component in the original coefficients and also
under each stated variable frame change.

\subsection{The algebraic limit and its two exact metric presentations}

Consider the algebraic direct limit of the injective transitions
\(V_{\ell m}\). A class represented by \(r\in\mathscr R_m^N\)
is sent to the original polynomial with coefficients \(L_m^Nr\).
Equation (NN.8) proves that the map is well-defined. It is onto
\(\mathcal P_\infty=\mathbb C[s_1,\ldots,s_k]\), since every
polynomial lies in a finite source. If its polynomial is zero then
\(r=\boldsymbol N_m L_m^Nr=0\), so it is injective. We have proved
the exact complex-linear isomorphism
\[
 \underset{m,V_{\ell m}}{\operatorname{colim}}\mathscr R_m^N
       \xrightarrow{\ \sim\ }\mathcal P_\infty.
 \tag{NN.10}
\]
Because (NN.7) is contractive, the limit seminorm is well-defined by
\[
 \|[r]\|_{N,\infty}
       =\lim_{\ell\to\infty}\|V_{\ell m}r\|
       =\inf_{\ell\ge m}\|V_{\ell m}r\|.
 \tag{NN.11}
\]
Representing the same class at a later index leaves the limit unchanged;
representing two classes at a common index proves homogeneity and the
triangle inequality. For the fixed polynomial given by (NN.10), (NN.4)
proves
\[
       \|[r]\|_{N,\infty}=0\quad\text{for every class}.
 \tag{NN.12}
\]
Thus its Hausdorff seminorm quotient, and its completion, are the zero
vector space. This conclusion retains the nonzero algebraic space in
(NN.10), together with the exact zero quotient map.

For comparison retain the complete derivative graph
\(\mathscr G_m=\operatorname{ran}(\boldsymbol T_m,\boldsymbol N_m)\)
with its original sum norm. Its transitions are
\[
 W_{\ell m}(\boldsymbol T_m p,\boldsymbol N_m p)
       =(\boldsymbol T_\ell Q_{\ell m}p,
                            \boldsymbol N_\ell Q_{\ell m}p).
 \tag{NN.13}
\]
Uniqueness of the source proves that these maps are well-defined,
injective and compose. Orthogonality and (NN.3) show that both sides
have squared norm \(\|\partial_u A_m p\|^2\), so every transition
is isometric. Addition of the two components commutes with transitions
and gives the exact isometry
\[
 \underset{m,W_{\ell m}}{\operatorname{colim}}\mathscr G_m
    \xrightarrow{\ \operatorname{Add}\ }
    \{\partial_u(\Psi P):P\in\mathcal P_\infty\}\subset\mathscr K.
 \tag{NN.14}
\]
Its extension identifies the Hilbert completion with the closure of
this specified derivative image in the original \(\mathscr K\).
The derivative image is injective in its polynomial source. Indeed
\(\Psi P\in L^2(du;\mathscr H)\) and a zero weak derivative in
\(u\) make every scalar pairing with an \(\mathscr H\) vector a
constant distribution in \(u\). Such a constant belongs to
\(L^2(du)\) only when zero. Hence \(\Psi P=0\). The nonzero
entire one-factor amplitude has a discrete real zero set, so \(\Psi\)
is nonzero almost everywhere in the original coordinates; a polynomial
vanishing almost everywhere there is the zero polynomial.

The algebraic reconstruction maps of (FC.31) give the commuting identity
\[
 W_{\ell m}\mathsf{Rec}_m
       =\mathsf{Rec}_\ell V_{\ell m}.
 \tag{NN.15}
\]
It follows by substituting \(r=\boldsymbol N_m p\). Thus the two
algebraic limits are explicitly isomorphic through the original source,
while (NN.12) and (NN.14) compute their two distinct retained metrics
and their exact completion maps.

\subsection{Original arithmetic jets and the size of reconstruction}

For a nonempty packet retain the original finite algebra
\(E=\mathbb C[\mathbf s]/(h(s_1),\ldots,h(s_k))\), its complete
Hermite coordinates, and \(U=\prod_i v_h(s_i)\). Its jet is a unit.
The fixed polynomial observation is the actual derivative jet
\[
 \mathcal J^{\log}P
    =j_I\left(\frac1k\sum_i\partial_{s_i}(UP)\right)\in E.
 \tag{NN.16}
\]
It includes the entire product derivative of \(U\) at every retained
nilpotent order. Formula (FC.34) gives its exact finite matrices
\(\mathsf J_m^{\log}\), including the original Mellin factors and
the \(-i\mathscr U_k^{-1}\) arrow from the derivative amplitude.
Their coefficient inclusions give
\(\mathsf J_\ell^{\log}Q_{\ell m}=\mathsf J_m^{\log}\), whence
\[
 \mathsf{Obs}_\ell^{\log}V_{\ell m}
       =\mathsf{Obs}_m^{\log},\qquad
 \mathsf{Obs}_m^{\log}=\mathsf J_m^{\log}L_m^N.
 \tag{NN.17}
\]
These equalities follow from (NN.8), so they define on (NN.10) exactly
the original algebraic map \(\mathcal J^{\log}\).

This map is nonzero, with an explicit original source witness. Choose
any retained tuple \(\boldsymbol\rho=(\rho_1,\ldots,\rho_k)\), set
\(\sigma=\sum_i\rho_i\), and take \(P_\rho=S-\sigma\).
The complete jet evaluation at the constant coordinate of that tuple is
\[
 (\mathcal J^{\log}P_\rho)(\boldsymbol\rho)
       =U(\boldsymbol\rho)\ne0.
 \tag{NN.18}
\]
The product-rule term multiplying \(P_\rho\) vanishes at that tuple;
the other term is \(U\partial_S^{\rm rel}S=U\), with its full
\(1/k\) retained. The nonvanishing follows from taking every selected
zero to its full order in \(h\).

There is also a permutation-invariant witness that is itself an
original theta relation. Take \(P_{\rm rel}=\sum_i h(s_i)\in I\).
Its original full jet is zero, while its retained derivative is
\[
 \mathcal J^{\log}P_{\rm rel}
       =\frac1k U_k\sum_i[h'(s_i)]_I\ne0.
 \tag{NN.18a}
\]
The term differentiating \(U\) still has the factor \(P_{\rm rel}\)
and is zero in \(E\). If \(d=1\), the sum of derivatives is exactly
\(k\). If \(d\ge2\), its degree-\(d-1\) standard remainder terms
are the distinct monomials \(d s_i^{d-1}\); their nonzero coefficients
cannot cancel. Multiplication by the full unit \(U_k\) proves the
nonvanishing in both cases, including all confluent orders. Each
\(h(s_i)\) term has the original theta primitive with the two matching
Koszul signs, so this witness retains the actual relation before its
supported-zero quotient. It is available in the full invariant source
as well as in the full tensor source.

Fix a genuine original target metric as follows. Choose a degree
\(m_*\ge k(d-1)\). The original jet map
\(J_*p=j_I(UP_p)\) on \(\mathcal P_{m_*}\) is onto: full monic
remainders have total degree at most \(k(d-1)\), and multiplication
by \(j_IU\) is invertible. With its original amplitude metric
\(H_{m_*}^0=A_{m_*}^*A_{m_*}\), the retained minimum-source metric is
\[
 G_*=(J_*(H_{m_*}^0)^{-1}J_*^*)^{-1}>0.
 \tag{NN.19}
\]
Indeed for a target vector \(q\), the source
\((H_{m_*}^0)^{-1}J_*^*G_*q\) maps to \(q\); it is orthogonal
in the original source metric to \(\ker J_*\). Its squared norm
is \(q^*G_*q\), and adding any kernel vector adds its nonnegative
squared norm. This proves the stated actual minimum and its constants.

The exact sizes of the finite inverse, graph reconstruction and
arithmetic observation are
\[
 \begin{aligned}
 \|L_m^N\|_{N\to0}^2
   &=\lambda_{\max}((H_m^N)^{-1/2}H_m^0(H_m^N)^{-1/2}),\\
 \|\mathsf{Rec}_m\|^2
   &=\lambda_{\max}((H_m^N)^{-1/2}H_m^\partial(H_m^N)^{-1/2}),\\
 \|\mathsf{Obs}_m^{\log}\|_{N\to G_*}^2
   &=\lambda_{\max}((H_m^N)^{-1/2}
          (\mathsf J_m^{\log})^*G_*\mathsf J_m^{\log}
                                            (H_m^N)^{-1/2}).
 \end{aligned}
 \tag{NN.20}
\]
Here \(H_m^0\), \(H_m^N\), and \(H_m^\partial\) are the three
original forms in (FC.29)--(FC.32). For a nonzero coefficient vector
\(p\), the denominator of each norm ratio is \(p^*H_m^Np\);
the respective numerators are \(p^*H_m^0p\),
\(p^*H_m^\partial p\), and
\(p^*(\mathsf J_m^{\log})^*G_*\mathsf J_m^{\log}p\).
The change of variable \(b=(H_m^N)^{1/2}p\), with its exact inverse,
proves the displayed eigenvalue formulas without altering any of these
original forms.

All three norms tend to infinity as \(m\to\infty\). For the first
two use any fixed nonzero polynomial: its original amplitude and its
derivative amplitude both have fixed positive norm, whereas (NN.4)
makes its normal norm tend to zero and (FC.26) keeps that norm strictly
positive at every finite degree. For the third use the single fixed
polynomial \(P_\rho\) in (NN.18); its \(G_*\) target norm is a
fixed positive number. Each operator norm is bounded below by this
specified ratio for every sufficiently large \(m\), proving divergence.
Equivalently, (NN.17) is a nonzero algebraic observation on the limit,
and cannot descend continuously through its zero Hausdorff quotient.

Every changing actual minimum metric is retained as well. Define
\(G_m\) by the same minimum-source construction at degree \(m\ge m_*\).
Its original-source inclusion makes \(G_m\preceq G_*\): the earlier
minimum lift is an allowed later lift with exactly the same amplitude
norm. Set \(B_m=G_*^{-1/2}G_mG_*^{-1/2}>0\). The precise comparison is
\[
 \begin{gathered}
 q^*G_m q=(G_*^{1/2}q)^*B_m(G_*^{1/2}q),\qquad 0<B_m\preceq I,\\
 \|\mathsf{Obs}_m^{\log}\|_{N\to G_m}^2
  =\lambda_{\max}((H_m^N)^{-1/2}
       (\mathsf J_m^{\log})^*G_m\mathsf J_m^{\log}(H_m^N)^{-1/2}),\\
 \lambda_{\min}(B_m)\|\mathsf{Obs}_m^{\log}\|_{N\to G_*}^2
 \le\|\mathsf{Obs}_m^{\log}\|_{N\to G_m}^2
 \le\lambda_{\max}(B_m)\|\mathsf{Obs}_m^{\log}\|_{N\to G_*}^2.
 \end{gathered}
 \tag{NN.21}
\]
The first equality is direct multiplication. Its minimum and maximum
eigenvalue inequalities applied to \(q=\mathsf J_m^{\log}p\), followed
by the same Rayleigh supremum, prove every remaining assertion. Thus
the divergence proved for the fixed actual metric is transported to
the changing actual metric by an explicit positive operator; its
displayed eigenvalues have not been assigned an unproved asymptotic.

\subsection{The split lift retains the algebraic source of the zero limit}

For a complex vector space \(V\), use the split semimodule
\(G(V)=\{\tau\}\sqcup V\) over \(G(\mathbb C)\). The external
\(\tau\) is the additive identity, addition of supported vectors is
their original vector addition, supported scalar multiplication is
the original one, and multiplication involving external \(\tau\)
is external \(\tau\). Its supported zero is \(e_V=0_V\ne\tau\).
For a linear map \(f:V\to W\), the map \(G(f)\) sends \(\tau\)
to \(\tau\) and a supported vector \(v\) to the supported vector
\(f(v)\). Checking separately external and supported inputs proves
additivity and compatibility with every scalar. Composition is the
original composition on supported inputs and fixes \(\tau\).

Consequently (NN.10), (NN.15), and (NN.17) give exact split-semimodule
maps. In particular the seminorm quotient in (NN.12) has the split lift
\[
 G(\mathcal P_\infty)\longrightarrow G(0)=\{\tau,e_0\},\qquad
 \tau\longmapsto\tau,\quad P\longmapsto e_0
                          \quad(P\in\mathcal P_\infty).
 \tag{NN.22}
\]
Its two fibres are exactly \(\{\tau\}\) and the entire supported
polynomial source. The algebraic arithmetic map is simultaneously
\(G(\mathcal J^{\log}):G(\mathcal P_\infty)\to G(E)\).
For a nonempty packet it cannot factor through (NN.22): evaluating
a proposed factorization on the supported zero polynomial forces the
image of \(e_0\) to be \(e_E\), whereas (NN.18) supplies a polynomial
whose supported image has nonzero amplitude. This is an exact
factorization obstruction between the explicitly constructed maps,
with the original source, supported zeros and external absence retained.

For \(h=1\), the algebraic target \(E\) is zero, so its arithmetic
observation is the zero map and its split lift does factor through
(NN.22). The amplitude, density theorem, injective finite normal maps,
inverse/graph divergence and the complete graph limit remain unchanged.
The arithmetic norm formulas involving a nonzero \(G_*\) target are
then replaced by their unique zero-target maps. For a full invariant
source under a fixed finite permutation group, all constructions
restrict through the exact Reynolds maps of the density section; the
derivative and source are invariant, \(P_\rho=S-\sigma\) is invariant,
and the arithmetic target is the corresponding invariant subspace with
the restriction of its original target form. Thus the same proofs
retain that invariant source and its original metric rather than adding
noninvariant vectors.
To verify the asserted minimum metric on this subspace, permutations
act unitarily on the original tensor amplitude space because they
preserve its product measure and its identical one-factor amplitudes.
The full jet map is equivariant. Averaging an arbitrary lift of an
invariant target with the literal factor \(1/|\mathfrak G|\) gives an
invariant lift of the same target, so the invariant jet map is onto.
Uniqueness of the full minimum shows that every permutation fixes that
minimum lift. It is consequently the minimum in the invariant source
as well, with exactly the restriction of \(G_m\) or \(G_*\).
